\subsection{Regularity for a measure with a H\"older reference}

The same smoothing argument applies to a singular solution measure dominated by the product of a H\"older potential. Its square-capacity bound supplies stability, and the weighted Chern estimate supplies the smoothing error.

\begin{theorem}\label{thm:holder-dominated-continuous-regularity}
Let $u\in\PSH(X,\omega)\cap C^{0,\alpha}(X)$, $0<\alpha\leq1$, and suppose $\varphi\in\PSH(X,\omega)\cap C(X)$ satisfies
\[
 \omega_\varphi^n=\mu\leq A\omega_u^n,
 \qquad A>0.
\]
Then $\varphi$ is H\"older continuous with exponent
\[
 \beta=\frac{\alpha}{(\alpha+2n)(6n+1)+\alpha}>0.
\]
Its H\"older constant is uniform for fixed bounds on $A$, $\|u\|_{C^{0,\alpha}}$ and the oscillation of $\varphi$.
\end{theorem}
\begin{proof}
Normalize $\sup_X\varphi=0$ and write $M=\|\varphi\|_\infty$. By \cref{prop:holder-reference-capacity}, the measure $\mu$ has a square-capacity bound with a fixed constant. Put
\[
 \eta=\frac{2\alpha}{\alpha+2n},\qquad
 \gamma=\frac1{6n+1},\qquad q=\eta\gamma/2<1.
\]
The weighted estimate of \cref{lem:holder-reference-chern-error} gives $\|R_\delta\varphi-\varphi\|_{L^1(\mu)}\leq C\delta^\eta$.

Use the constants $A_0,K$ of \cref{lem:chern-kiselman-input}, taking $A_0\geq1$, and define the unshifted transform
\[
 F=\inf_{0<t\leq\delta}
       \{R_t\varphi+Kt^2+Kt-b\log(t/\delta)\}.
\]
It differs from $\Phi_{b,\delta}$ by the constant $K(\delta^2+\delta)$, so
\[
 \omega+\ddc F\geq-(A_0b+2K\delta)\omega,
 \qquad \varphi\leq F\leq R_\delta\varphi+2K\delta
\]
when $\delta\leq1$. Choose $\varepsilon=\delta^q$ and $b=(\varepsilon-2K\delta)/A_0$. For small $\delta$, $\varepsilon\leq1/2$ and $\varepsilon/(2A_0)\leq b\leq\varepsilon/A_0$.

Set $T=(1-\varepsilon)F$. Then $\omega+\ddc T\geq\varepsilon^2\omega$. The transform is continuous: its positive logarithmic penalty excludes radii near zero uniformly in the base point for each fixed $\delta,b$, and the minimum of its continuous defining function on the resulting compact interval depends continuously on the base point. Take $B\geq M+2K+1$ and set $\psi=T-B\varepsilon$. Since $R_\delta\varphi\leq0$, we have $-M-B\leq\psi\leq0$. Moreover,
\[
 \psi-\varphi
 \leq(1-\varepsilon)(R_\delta\varphi-\varphi+2K\delta)
                         +\varepsilon(M-B)
 \leq(1-\varepsilon)(R_\delta\varphi-\varphi+2K\delta).
\]
The weighted smoothing estimate and the fixed bound for $\mu(X)$ give $\int(\psi-\varphi)_+\,\mathrm d\mu\leq C\delta^\eta$. Thus \cref{prop:capacity-square-one-sided-stability}, with a fixed competitor bound $H\geq M+B+1$, gives
\[
 \sup_X(T-\varphi)\leq C\delta^{\eta\gamma}+B\varepsilon\leq C_1\varepsilon.
\]
Since $\varphi\leq0$,
\[
 0\leq F-\varphi
 =\frac{T-\varphi+\varepsilon\varphi}{1-\varepsilon}
 \leq2C_1\varepsilon.
\]
At a minimizing radius $t_0\in(0,\delta]$, monotonicity of $R_t\varphi+Kt^2$ implies
\[
 b\log(\delta/t_0)\leq F-\varphi\leq2C_1\varepsilon,
 \qquad t_0\geq\kappa\delta,
 \quad\kappa=\exp(-4A_0C_1)>0.
\]
The constant $\kappa$ is independent of $\delta$ and the point. Applying monotonicity again gives $R_{\kappa\delta}\varphi-\varphi\leq2C_1\delta^q$. With the lower bound $R_r\varphi-\varphi\geq-Kr^2$ and $r=\kappa\delta$, this proves $\|R_r\varphi-\varphi\|_\infty\leq Cr^q$.

Finally \cref{lem:chern-smoothing-integral} gives $\|\nabla R_r\varphi\|_\infty\leq C/r$. For points a distance $s$ apart,
\[
 |\varphi(x)-\varphi(y)|\leq2Cr^q+Cs/r.
\]
Take $r=s^{1/(1+q)}$ for small $s$; boundedness handles larger distances. The exponent is $q/(1+q)=\beta$.
\end{proof}
